\section{Background and Threat Model}
\label{sec:background}

\subsection{MCP tools and what a client can see}

An MCP server lists its tools through \texttt{tools/list}. Each declaration
contains a name, a natural-language description, and a JSON schema for the
input. A client invokes a tool through \texttt{tools/call} with a name and
arguments, and the server returns content together with an optional error
flag~\cite{mcp2025tools}. Declarations may also carry annotations such as
\texttt{readOnlyHint} or \texttt{destructiveHint}. The specification itself
states that these annotations should be treated as untrusted unless the server
is trusted. Nothing in the protocol gives the client an independent view of
what the call changed.

It helps to write this down. Let a tool implementation be a function $f$ that
maps arguments $a$ to a response $r$ and an effect $e$ on some state:
$(r, e) = f(a)$. An ordinary client observes $a$, because it sent it, and $r$,
because the server returned it. It does not observe $e$. Everything a response
auditor concludes about $e$ is therefore inferred from quantities the server
controls.

\subsection{Deployment model}

We consider a client that runs MCP servers locally, which is the common setup
for desktop agents and coding assistants. The client starts each server as a
process, speaks to it over standard input and output, and gives it access to
some part of the machine, such as a project directory or a database file. The
client, the operating system, and a small trusted mediator run with the user's
authority. Server packages come from public registries and are updated by their
authors.

\subsection{Adversary}

The adversary controls the server implementation after the user has approved
it. This models a malicious update of a package that was honest when it was
reviewed, a compromised dependency, or a server that behaves well until it sees
a particular request. The adversary may keep the declaration unchanged. It may
return the honest response while performing a different effect. It may write a
different path, different bytes, extra files or rows, links, or nothing at all.
It may start child or background processes, race the mediator's checks, replay
an approved request, and answer any follow-up read with values chosen to look
consistent. It knows how MCPGate works.

The adversary cannot modify the client, the mediator, the trusted store, or the
operating system's isolation. It cannot write the trusted store directly. It
cannot forge a contract.

We trust the contract as an input. MCPGate binds a server to what was
approved; it does not decide whether what was approved matches the user's
intent. Indirect prompt injection, where untrusted content steers the agent
into requesting a harmful call~\cite{greshake2023injection,zhan2024injecagent},
is therefore a separate problem that is handled before a contract exists. The
prototype also does not implement contract signing or expiry. It gives each
contract a canonical, immutable identity and treats authenticity and freshness
as a deployment assumption.

\subsection{Security goal and non-goals}

The property MCPGate provides is \emph{trusted-state admission}.

\begin{quote}
If MCPGate reports that a call committed, then the objects written to the
trusted store are exactly the objects that the mediator read from the frozen
private result, and that result satisfied the call's contract.
\end{quote}

For an exact contract this means the committed path and bytes equal the
approved path and bytes. For a contract that constrains only part of the
content, it means the committed objects satisfy every rule the contract states,
and nothing else. We report which content is constrained and which is not, so
a reader can see exactly how strong the guarantee is for each workflow.

The property is not \emph{world-effect prevention}. If the server can reach a
channel that the private copy does not contain, such as an open network socket
or a path outside its sandbox, it can act there before MCPGate refuses the
result. MCPGate does not protect confidentiality, does not prevent denial of
service, and does not roll back effects outside the mediated state. These
concerns belong to the confinement layer beneath MCPGate, and we compose with
such a layer rather than replace it.

We also separate two outcomes that are easy to conflate. \emph{Confinement}
means no unapproved effect entered trusted state. \emph{Completion} means the
approved effect did enter it. A server that silently does nothing is confined
but not completed. When MCPGate cannot establish that every writer has stopped,
or cannot observe the final state, it refuses or reports UNKNOWN rather than
guessing.
